% CS F407 Artificial Intelligence: final project report template.
% Limit: 2 pages, including everything except references.
% Do not change the font size, margins or spacing.
\documentclass[11pt,a4paper,twocolumn]{article}

\usepackage[margin=2cm,columnsep=0.8cm]{geometry}
\usepackage[T1]{fontenc}
\usepackage{lmodern}
\usepackage{microtype}
\usepackage{amsmath,amssymb}
\usepackage{graphicx}
\usepackage{booktabs}
\usepackage{tabularx}
\usepackage{enumitem}
\usepackage{titlesec}
\usepackage[hidelinks]{hyperref}

\titleformat{\section}{\normalsize\bfseries}{\thesection}{0.6em}{}
\titlespacing*{\section}{0pt}{1.2ex plus 0.3ex}{0.6ex}
\setlist{nosep,leftmargin=*}
\setlength{\parskip}{0.3ex}

\newcommand{\coursename}{CS F407 Artificial Intelligence, AY 2026--27 Semester I}
\newcommand{\teamno}{Team CTRL+AI}
\newcommand{\githuburl}{https://github.com/krishavgoyal/Cross_Era_Fantasy_Cricket}

\begin{document}

\twocolumn[{%
\begin{center}
  {\large\bfseries Cross-Era Fantasy Cricket: Animating Historical Cricketers with Modern Motion}\\[1ex]
  \teamno\\
  Pranahv Subramaniam Sankar (2024A4PS0423G), Ahaan Pathak (2024A3PS0396G),\\
  Ankit Pandey (2024AAPS0808G), Krishav Goyal (2024A7PS0855G)\\[0.5ex]
  {\small \coursename, BITS Pilani, K.K. Birla Goa Campus}\\[1.5ex]
  \begin{minipage}{0.85\textwidth}
    \small
    \centerline{\bfseries Abstract}\vspace{0.5ex}
    Fantasy cricket compares players across eras only through statistics. We make such match-ups visible: from
    one archival photograph and a modern clip we generate a short video of the historical player performing the
    modern motion, styled like footage of his era. DWPose extracts the motion on a laptop CPU; Wan2.1-VACE-1.3B
    generates the clip on a free Kaggle T4 from the photo and a control video; post-processing retimes it and
    applies 1930s/1970s/1990s filters. In a controlled ablation, a depth control instead of a pose skeleton kept
    the bat in 5/5 runs (0/2 with the skeleton) and lowered brightness flicker from 2.38 to 1.68, with motion
    accuracy (PCK@0.2) of 0.90--1.00. Our final video shows Harold Larwood bowling to Sir Don Bradman, and a
    catalog renders the match in three era styles.
  \end{minipage}
\end{center}
\vspace{1.5ex}
}]

\section{Introduction}
What if Harold Larwood bowled to Sir Don Bradman, using motion filmed today? \textbf{Input:} one photograph of
a historical player (appearance) and a clip of a modern player (motion). \textbf{Output:} a 2--6\,s video of
the historical player performing that motion in the visual style of his era, labelled ``AI-generated''. We
contribute (i)~a free pipeline (laptop CPU, free Kaggle GPU); (ii)~a player-only depth control that keeps the
bat and removes bystanders; (iii)~a multi-seed ablation; (iv)~a cross-era match video and catalog; and (v)~an
analysis of why a skeleton-driven baseline fails on fast bowling.

\section{Background and related work}
This is reference-guided human image animation. DWPose~\cite{dwpose} gives 2D joint sequences. Diffusion
models with a reference branch (Animate Anyone~\cite{animateanyone}, MagicAnimate~\cite{magicanimate},
MusePose~\cite{musepose}) animate a photo from poses; AnimateDiff~\cite{animatediff} adds temporal layers to
image diffusion. Video diffusion
transformers such as Wan2.1~\cite{wan} with VACE~\cite{vace} accept a reference image and a pose or depth video.
Depth Anything V2~\cite{dav2} estimates monocular depth. The proposal named LTX-2~\cite{ltx2}, which uses the
photo only as a first frame and needs 32\,GB+ GPUs, and MagicAnimate (2023), kept as a team baseline; we
chose VACE.

\section{Method}
\textbf{Motion.} rtmlib DWPose-l runs on CPU; the player is chosen by a box on frame 0 and followed by best-IoU
tracking. Joints with confidence $<0.3$ are masked, gaps of up to 4 frames interpolated, tracks smoothed
(Savitzky--Golay, window 7) and resampled to 33 frames at 16\,fps on a 480$\times$832 canvas with an 8\% top
margin. \textbf{Controls.} \emph{Skeleton}: an OpenPose-style render. \emph{Depth}: Depth Anything V2 on the
same frames and crop, globally normalised (no flicker), head flattened so a modern helmet does not replace the
face. \emph{Player-only depth} (bowler): a person mask kept only around the player's keypoints, so fielders
vanish. \textbf{Reference.} rembg removes the photo background; the
player is tightly cropped on white. \textbf{Generation.} Wan2.1-VACE-1.3B (fp16) in a headless ComfyUI on a
Kaggle T4: 20 steps, CFG 5, shift 8. The negative prompt lists observed failures (extra balls, two bats,
logos, hats). \textbf{Post-processing.} Frame blending retimes the batter to real speed
at 24\,fps; the bowler is shown in 2$\times$ slow motion with frame holds, since blending ghosted the fast arm.
Seeded era filters (1930s: grain, flicker, gate weave, 18\,fps; 1970s: scanlines, colour fringing, 25\,fps;
1990s: chroma bleed, 25\,fps) burn in the label.

\section{Experimental setup}
\textbf{Data.} A modern cover drive (5.6\,s) and a front-on fast bowler (2.0\,s) supply motion; a portrait of
Bradman and a public-domain 1929 portrait of Harold Larwood (National Library of Australia) supply appearance.
\textbf{Runs.} A 2$\times$2 ablation, control (skeleton/depth) $\times$ reference (padded/tight crop), 17
frames; seeds 42/7/123/2024 for depth with the tight crop; final 33-frame runs: batter F1 (seed 42) and bowler
L1--L3 (seeds 42/7/2024). \textbf{Metrics.} PCK@0.2: DWPose on the output vs the control joints (17 joints,
threshold 0.2$\times$ torso length). Identity: DINOv2~\cite{dinov2} cosine similarity of person crops vs the
photo. Flicker: std of per-frame brightness. Bat: one rigid bat in every frame, by inspection.
\textbf{Hardware.} The laptop CPU for everything except generation; on the Kaggle T4 a 33-frame clip takes
${\approx}520$\,s at a 14.9\,GB peak, so clips stay at 33 frames. \textbf{Baselines.} Skeleton
control; the team's DensePose-driven \mbox{AnimateDiff}/MagicAnimate engine; a skeleton-driven MusePose
pipeline (Section~\ref{sec:disc}).

\section{Results}
\begin{table}[h]
\centering
\small
\begin{tabular}{lcccc}
\toprule
Control (runs) & PCK & ID & Flicker & Bat \\
\midrule
Skeleton (A, C) & 0.95 & 0.19 & 2.38 & 0/2 \\
Depth (B, D1--D4) & 0.97 & 0.23 & 1.68 & 5/5 \\
\midrule
Depth, batter F1 & 0.90 & 0.22 & 1.65 & yes \\
Bowler L1--L3 & 0.81 & 0.11 & 3.67 & n/a \\
Bowler L3 (chosen) & 0.80 & 0.12 & 1.81 & n/a \\
\bottomrule
\end{tabular}
\caption{Means over runs; bowlers use player-only depth. Higher PCK@0.2 and ID (identity), lower flicker are
better.}
\label{tab:results}
\end{table}

\textbf{Depth fixes the bat.} With only a skeleton, which never describes the bat, the model loses or bends it;
with depth, one rigid bat is held in every frame of all five runs, with no loss of motion accuracy and less
flicker (Table~\ref{tab:results}). \textbf{The tight crop did not help likeness}: with the seed fixed it changed
identity by $-0.02$ (vs the crop) and $-0.08$ (vs the original photo), while the seed alone moved it by up to
0.14, so likeness comparisons need several seeds. \textbf{Bowler.} Player-only depth removed the fielders, and
the negative prompt removed the photo's fedora, suit and cigarette. PCK is lower (0.79--0.85) because the
bowling arm is the fastest, most blurred motion; L3 was chosen for its period kit and lowest flicker.
\textbf{Eras and final video.} Era filters keep the motion: PCK falls by at most 0.08 (batter) and 0.12
(bowler), while the 1930s gate weave raises warping error (5.7 to 9.8 for the batter). The final video joins
Larwood (L3, slow motion) and Bradman (F1, real speed), both as 1930s film; a catalog renders the same match as
1930s film, 1970s TV and 1990s broadcast, plus a style-swap version. \textbf{Team baseline.} \mbox{AnimateDiff}/MagicAnimate gave 384$\times$512 clips of 16 keyframes at 7\,fps,
too coarse for a match.

\section{Discussion and limitations}\label{sec:disc}
\textbf{Why the skeleton baseline failed on bowling.} It fits a rigid 3D skeleton to DWPose and
MediaPipe~\cite{mediapipe} joints and drives MusePose. At a stricter PCK@0.1 (not comparable to
Table~\ref{tab:results}) it scored 0.963 on batting but 0.64 on bowling: its dark, side-on slow-motion bowling
clip blurred the arm and lost the feet, and a skeleton does not say which limb is in front, so the bowler was
detected mirrored in up to 60\% of frames (0.88 if mirroring is allowed). The VACE bowler used a sharp,
front-on clip and depth, which fixes front/back order. \textbf{Limitations.} Face likeness is
weak (identity 0.1--0.3), the helmet outline can become a dark cap, clips are short (33 frames), and the players
are generated separately. \textbf{Ethics.} Every frame is labelled AI-generated; modern clips are used only for motion, not
redistributed, with overlays blurred. \textbf{Next:} longer clips and a face adapter.

\section{Conclusion}
A free pipeline animates historical cricketers convincingly for batting. Depth control matters more than the
reference crop, and clean footage matters most for bowling.

\section*{Individual member contributions}
{\small
\begin{tabularx}{\linewidth}{@{}l X@{}}
\toprule
Name & Contribution \\
\midrule
Ankit Pandey & Problem formulation, pipeline overview, limitations and ethics, Sections 1, 2 and 7. \\
Pranahv S.\ Sankar & Team pose engine, rigid 3D skeleton fit and the MusePose baseline, Section 6. \\
Ahaan Pathak & Model survey and the AnimateDiff/MagicAnimate team baseline. \\
Krishav Goyal & VACE pipeline and depth controls, ablation, bowler, era filters, evaluation, catalog,
video, Sections 3--5. \\
\bottomrule
\end{tabularx}}

\smallskip
\noindent All members have read this report and agree with the statement above.

\section*{Generative AI disclosure}
We used Claude (Anthropic, Claude Opus 5.5) in the Claude app for research, guidance, drafting this report
and editing the video (captions, voice clean-up), and Claude Code (Claude Opus 5.5) to write and debug the
code in \texttt{scripts/}, \texttt{kaggle/}, \texttt{motion3d/} and \texttt{animate/}. We checked all output,
ran the code ourselves, verified every number against our runs, and take responsibility for it.

\section*{Code availability}
\noindent Code and data:\\
\url{\githuburl}\\
The repository is public and its README gives instructions to reproduce the results reported here.

\begin{thebibliography}{99}
\small
\bibitem{dwpose}
Z.~Yang, A.~Zeng, C.~Yuan and Y.~Li. Effective whole-body pose estimation with two-stages distillation.
arXiv:2307.15880, 2023.
\bibitem{animateanyone}
L.~Hu, X.~Gao, P.~Zhang, K.~Sun, B.~Zhang and L.~Bo. Animate Anyone: consistent and controllable image-to-video
synthesis for character animation. arXiv:2311.17117, 2023.
\bibitem{magicanimate}
Z.~Xu, J.~Zhang, J.~H.~Liew, H.~Yan, J.-W.~Liu, C.~Zhang, J.~Feng and M.~Z.~Shou. MagicAnimate: temporally
consistent human image animation using diffusion model. arXiv:2311.16498, 2023.
\bibitem{musepose}
Z.~Tong, C.~Li, Z.~Chen, B.~Wu and W.~Zhou. MusePose: a pose-driven image-to-video framework for virtual human
generation. \url{https://github.com/TMElyralab/MusePose}, 2024.
\bibitem{animatediff}
Y.~Guo, C.~Yang, A.~Rao, Z.~Liang, Y.~Wang, Y.~Qiao, M.~Agrawala, D.~Lin and B.~Dai. AnimateDiff: animate your
personalized text-to-image diffusion models without specific tuning. arXiv:2307.04725, 2023.
\bibitem{wan}
Wan Team. Wan: open and advanced large-scale video generative models. arXiv:2503.20314, 2025.
\bibitem{vace}
Z.~Jiang, Z.~Han, C.~Mao, J.~Zhang, Y.~Pan and Y.~Liu. VACE: all-in-one video creation and editing.
arXiv:2503.07598, 2025.
\bibitem{dav2}
L.~Yang, B.~Kang, Z.~Huang, Z.~Zhao, X.~Xu, J.~Feng and H.~Zhao. Depth Anything V2. arXiv:2406.09414, 2024.
\bibitem{ltx2}
Y.~HaCohen et al.\ (Lightricks). LTX-2: efficient joint audio-visual foundation model. arXiv:2601.03233, 2026.
\bibitem{dinov2}
M.~Oquab et al. DINOv2: learning robust visual features without supervision. arXiv:2304.07193, 2023.
\bibitem{mediapipe}
C.~Lugaresi et al. MediaPipe: a framework for building perception pipelines. arXiv:1906.08172, 2019.
\end{thebibliography}

\end{document}
